\section{Related Work}
\subsection{Occlusion in fruit perception: active versus passive viewpoints}
Occlusion is a long-standing bottleneck for agricultural robots. Bac \etal~\cite{bac2014} reviewed 50 harvesting robots and found average fruit localisation success of 85\% and harvest success of 66\%, and that performance had not improved over three decades. Occlusion was a frequent cause of failure. The dominant response has been active perception. In a systematic review of active perception for fruit harvesting, Magalhães \etal~\cite{magalhaes2022} screened 1,034 publications, kept 195, and found multiple-viewpoint strategies and visual servoing to be the most common, though not yet robust. Recent work plans the next best view of a manipulator-mounted camera to reveal task-relevant plant parts in occluded greenhouse crops~\cite{burusa2024}. Hemming \etal~\cite{hemming2014} quantified why this helps: combining camera positions raised sweet pepper detectability from at most 69\% to 90\%. All of these approaches choose a viewpoint per target, which assumes a manipulator and time spent at each fruit; a vehicle on a straight vine row cannot, and must use the views its motion happens to supply.

\subsection{Vineyard yield estimation and hidden fruit}
Nuske \etal~\cite{nuske2014} mounted a side-facing camera with flash illumination on a vehicle, detected berries in each image and calibrated counts to harvest yield over several seasons. Their estimates explained up to 75\% of spatial yield variance, with an average error of 3--11\% of total yield. Santos \etal~\cite{santos2020} segmented grape clusters with Mask R-CNN and associated detections across video frames in 3D to avoid double counting. Liu \etal~\cite{liu2019} counted orchard fruit with a single camera, tracking detections and building a semantic structure-from-motion map, and matched a system that used LiDAR and GPS/INS. These systems count what is visible and treat occlusion as a bias to correct afterwards.

Several methods model the hidden fruit directly. Millan \etal~\cite{millan2018} used a Boolean model to compensate statistically for partly occluded berries, obtaining $R^2=0.78$ for on-the-go yield on three vine segments. GCNet~\cite{quinones2025gcnet} is the closest to our setting. It images grape clusters indoors against white or simulated backgrounds with artificial foliage, segments them with a U-Net, and adds a regressed correction factor for occluded berries when estimating count and weight. Kierdorf \etal~\cite{kierdorf2022} go further and use a conditional GAN to generate a plausible view of berries behind leaves, which brought counts closer to manual references. These methods estimate hidden fruit from priors learned on other bunches. The result can be accurate on average, but for any given vine the hidden berries are inferred, not observed. Our method takes the opposite position: it recovers only fruit that some other view has actually seen, and leaves the rest empty. A learned correction could then be applied to a smaller remaining gap.

For berry-level counting, ViViD-5K~\cite{tong2026vivid} provides 5,000 vineyard images with over 648,000 berry centroids across 13 varieties, together with GrapeSAM, a point localiser followed by SAM. It reports a counting MAE of 21.46 berries per image on images averaging about 130 berries. We use its published point model unchanged as an independent counter. Its authors note weaker performance on indoor and uniform backgrounds, which applies to our scenes.

\subsection{Seeing through occluders with multiple views}
Synthetic aperture imaging aligns images from many camera positions on a chosen depth, so occluders off that plane blur away. Vaish \etal~\cite{vaish2006} compared focus, stereo and robust statistics such as per-pixel medians for reconstructing surfaces behind foliage. Airborne optical sectioning applies the same principle from a drone to reveal the ground through forest canopy~\cite{kurmi2018}. Each of these requires the target depth or a calibrated camera pose. Obstruction-free photography~\cite{xue2015} removes fences and reflections from a short handheld sequence by separating layers by their motion, and Liu \etal~\cite{liu2020obstruct} learn the same layer separation with optical flow and a CNN. DeOccNet~\cite{wang2020} learns occlusion removal from light fields. Learned methods produce visually complete images, but they inpaint regions no view has seen, which is undesirable when the output is to be counted. Our pipeline sits between these: like obstruction removal, it needs no depth or calibrated poses, but like robust synthetic aperture methods, it never synthesises content that no view observed. Pixel labelling of occluders~\cite{pei2013} and all-in-focus synthetic aperture~\cite{yang2014} are the published remedies for the occluder leakage and multi-depth blur of plain averaging (Sec.~\ref{sec:evolution}). We build on SIFT~\cite{lowe2004}, RANSAC~\cite{fischler1981} and SAM~3~\cite{carion2025}; Sec.~\ref{sec:method} gives the details.
